\begin{afterword}
\summaryhead

The introduction to the third book, \textsc{The Machine in the Ghost}, begins with the reader as a mind: a lawful product of evolution that can study its own workings and improve them for its own goals rather than evolution's. Our mental categories do not resemble physical ones. Philosophers once concluded that minds are separate from brains, the view Gilbert Ryle called ``the dogma of the Ghost in the Machine.'' Rejecting that view, the introduction says, has not by itself given us a predictive theory of the mind, so our talk of rationality may remain as unconstrained as talk of spirits. The book will therefore study processes that look purposeful but are not human, evolution and artificial intelligence, and the introduction previews its three sequences.

The second half introduces Yudkowsky's work on AI: the intelligence explosion, Friendly AI and existential risk, with a quotation from Russell and Norvig's textbook. It says these views are still debated. AI belongs in a book on human rationality, it argues, as a source of simple illustrations of cognition and because long-term forecasting is an important application of Bayesian rationality.

\responsehead

The introduction does two jobs. The first half sets up the book's theme, that the mind is a machine built by evolution and can be studied as one. The second half is a profile of the author's work on artificial intelligence. It reports its sources accurately: Ryle's phrase, and the quotation from ``The Martial Art of Rationality'' with small changes that do not alter the meaning. It also states the uncertainty around the author's views plainly. Little is known about when general AI might come, and forecasters have not reached a consensus.

The first half's main claim is a premise, not a finding. Rejecting dualism, it says, has not produced a better predictive model, and ``Practically speaking, our purposes and desires still function like free-floating ghosts.'' The claim is hedged, put as a conditional, and given no example. An introduction may leave the case to the essays that follow, but the reader should know that it is assumed here.

The case for including AI rests on two reasons. The first, that AI gives simple illustrations of cognition, is a plan for the book. The second rests on two sentences. People, experts included, are ``extraordinarily bad'' at forecasting major events; this has no source here, but Tetlock's forecasting tournaments support it. And Bayesian rationality ``can model correct reasoning even in domains where the data is scarce or equivocal.'' That is true in the sense that Bayes's rule says how to revise any starting estimate on the evidence available. With scarce data the result depends mostly on the starting estimate, which consistency does not fix. The two sentences together show that long-range forecasting is hard and that a consistent method exists, not that the method makes such forecasts reliable. The introduction does not claim that it does.

The context of the second half is left for the reader to supply. MIRI appears without explanation; it is the institute where Yudkowsky worked, and it released this book. The Russell and Norvig passage, in the textbook, opens as a report of Yudkowsky's own work. The introduction says the views are debated but cites no one who disputes them.

In short: an introduction that states its author's views on AI with their uncertainty, cites no one on the other side of the debate it mentions, and justifies AI's place in the book with a claim about Bayesian forecasting that holds only in a limited sense.
\end{afterword}
